\section{问题四：分区与独立资源配置}

\subsection{不可拆服务区组件和配置计算}

同一运输架次所访问的多个服务区必须同组。把15个服务区作为顶点，将同架次访问的服务区连边后取连通分量，得到四个不可拆组件。$C_1$包含S001、S002、S010、S012、S013、S014；$C_2$包含S003、S006、S007、S015；$C_3$包含S004、S005、S008、S011；$C_4$仅包含S009。拆分第三问中S002与S004的跨区航次，使原组件分成两部分。两组方案枚举七种非对称二分，三组方案枚举六种。图~\ref{fig:partition}给出所选两组与三组划分。

对固定时间表，某组所需同型实体数等于所属任务占用区间的最大并发数；所需同型共享电池数等于任务区间延长到电池充满时刻后的最大并发数。中继机区间延长给定周转时间，中继能源组件区间延长至充满时刻。同一问题三中继架次若保障多个任务组，则独立执行时各组均需复制等价架次和资源，不能跨组借用原架次。上述峰值法对固定区间任务是精确的最少资源数。

\begin{figure}[H]\centering\includegraphics[width=.94\textwidth]{../figures/partitions.pdf}\caption{保持第三问航次不变的两组与三组服务区划分}\label{fig:partition}\end{figure}

\subsection{两组与三组比较}

配置向量的八个分量依次表示A、B、C型运输机，A、B、C型电池，中继机和中继能源组件。按“库存缺口合计、配置数量、组间工作量变异系数”选出的两组方案为$C_1\cup C_2\cup C_3$与$C_4$，分别服务77箱和3箱。三组方案为$C_1\cup C_3$、$C_2$、$C_4$，分别服务55箱、22箱和3箱。各组配置、合计需求与库存缺口见表~\ref{tab:q4}。

\begin{table}[H]\centering\small\caption{独立分区的资源需求与缺口}\label{tab:q4}
\begin{tabular}{lcccccccc}\toprule
方案 & A机 & B机 & C机 & A电 & B电 & C电 & 中继机 & 中继件\\\midrule
库存 & 4&2&2&6&4&4&2&6\\
两组第1组 &4&2&2&6&4&4&2&4\\
两组第2组 &0&1&0&0&1&0&2&2\\
两组合计 &4&3&2&6&5&4&4&6\\
两组缺口 &0&1&0&0&1&0&2&0\\
三组第1组 &4&2&2&5&4&2&2&4\\
三组第2组 &1&1&2&1&2&2&1&1\\
三组第3组 &0&1&0&0&1&0&2&2\\
三组合计 &5&4&4&6&7&4&5&7\\
三组缺口 &1&2&2&0&3&0&3&1\\
\bottomrule\end{tabular}\end{table}

两组工作时长变异系数为0.9379，三组为0.9056；所选两组方案把77箱与3箱分开，均衡性较差，这是优先压缩库存缺口的代价。三组相对稍均衡，却要复制更多同时段占用的实体和中继资源。两种严格独立执行方式均超出现有库存，因此不能把第三问的全局共享资源可行性直接转移到分区方案。若不增加库存，需要放宽“任务安排及通信关系不变”的限制并重新联合排程，此处没有擅自改动题设。
